\chapter{Elliptic Curves}

\chapterSubhead{Cryptography's quiet workhorse---and another casualty of Shor}

If RSA is the famous lock, \keyterm{elliptic-curve cryptography} (ECC) is the small, efficient one that does most of the actual work. Introduced independently by \citet{koblitz_elliptic_1987} and \citet{miller_use_1986}, ECC provides the same security as RSA at much smaller key sizes, and has largely displaced RSA wherever bandwidth and CPU matter: TLS in modern browsers, mobile authentication, Bitcoin and most other cryptocurrencies, modern signing protocols (Ed25519, ECDSA). And, like RSA, it falls to Shor's algorithm.

%------------------------------------------------------------------
\section{What an elliptic curve is}
%------------------------------------------------------------------

An \keyterm{elliptic curve} over a field $\mathbb{F}$ is the set of points $(x,y) \in \mathbb{F}^2$ satisfying a Weierstrass equation
\[
  y^2 = x^3 + ax + b,
\]
together with a special point at infinity $\mathcal{O}$. For cryptography, we work over a finite field $\mathbb{F}_p$ (the integers mod a large prime $p$) or over a binary field $\mathbb{F}_{2^m}$. The condition $4a^3 + 27b^2 \neq 0$ (non-vanishing discriminant) ensures the curve is non-singular.

The points of an elliptic curve form an abelian group under a geometric ``chord-and-tangent'' addition law. Given two points $P, Q$ on the curve:

\begin{itemize}
  \item Draw the line through $P$ and $Q$ (or the tangent at $P$ if $P = Q$).
  \item Find the third intersection point $R'$ of this line with the curve.
  \item Reflect $R'$ across the $x$-axis. The reflected point is $P + Q$.
\end{itemize}

This rule, made explicit in coordinates, gives algebraic formulas for $P+Q$ and $2P$ in terms of the coordinates of $P$ and $Q$. The point $\mathcal{O}$ acts as the identity, and $-P$ is the reflection of $P$ across the $x$-axis. The group law is commutative and associative, making the points of the curve into an abelian group.

\begin{example}
On the curve $y^2 = x^3 - x + 1$ over $\mathbb{R}$, take $P = (0,1)$ and $Q = (1,1)$. The line through them is $y = 1$. Substituting: $1 = x^3 - x + 1$, so $x^3 - x = x(x-1)(x+1) = 0$. The third intersection is $R' = (-1, 1)$. Reflecting: $R = P + Q = (-1, -1)$.
\end{example}

%------------------------------------------------------------------
\section{Scalar multiplication and the discrete log problem}
%------------------------------------------------------------------

For an integer $k$ and a point $P$ on the curve, \keyterm{scalar multiplication} is
\[
  [k]P = \underbrace{P + P + \cdots + P}_{k \text{ times}}.
\]
This is fast to compute: by the standard ``double-and-add'' procedure, $[k]P$ requires $O(\log k)$ point doublings and additions, each of which is fast over $\mathbb{F}_p$.

The reverse problem, the \keyterm{elliptic curve discrete logarithm problem} (ECDLP), is hard:

\textbf{Problem.} Given a point $P$ on the curve and a point $Q = [k]P$ for some unknown integer $k$, recover $k$.

The fastest classical algorithm for the ECDLP on a well-chosen curve is Pollard rho, with complexity $O(\sqrt{n})$ where $n$ is the order of the subgroup generated by $P$. For a 256-bit curve order, this is $2^{128}$ operations---infeasible on any classical computer, now or in the foreseeable future.

The hardness of the ECDLP is, by every classical metric, more robust than the hardness of factoring. Factoring has the number-field sieve (sub-exponential); ECDLP on standard curves has only Pollard rho (fully exponential in the bit-length). This is why ECC keys can be much shorter than RSA keys for the same security level: a 256-bit ECC key delivers $\sim 128$ bits of security, comparable to a 3072-bit RSA key.

%------------------------------------------------------------------
\section{ECDH key exchange and ECDSA signatures}
%------------------------------------------------------------------

The two dominant ECC applications are key exchange and digital signatures.

\textbf{ECDH (Elliptic Curve Diffie--Hellman).}
\begin{enumerate}
  \item Alice and Bob agree on a curve $E$ over $\mathbb{F}_p$ and a base point $G$.
  \item Alice picks random private key $a$, computes public key $A = [a]G$.
  \item Bob picks random private key $b$, computes public key $B = [b]G$.
  \item Each exchanges their public key over an insecure channel.
  \item Both compute $[a][b]G = [ab]G$. This is their shared secret.
\end{enumerate}

An eavesdropper sees $A$ and $B$ but cannot extract $a$ or $b$ without solving the ECDLP. ECDH is the basis of the key exchange in modern TLS.

\textbf{ECDSA (Elliptic Curve Digital Signature Algorithm).} A digital signature scheme. To sign a message $m$ with private key $a$:
\begin{enumerate}
  \item Choose random nonce $k$, compute $R = [k]G = (x_R, y_R)$. Set $r = x_R \bmod n$ (where $n$ is the order of $G$).
  \item Compute $s = k^{-1}(h(m) + a r) \bmod n$, where $h$ is a hash function.
  \item Signature is $(r, s)$.
\end{enumerate}
Verification recovers $r$ from the signature and the public key $A$. A unique-nonce flaw (re-using $k$ for two different messages) allows recovery of the private key; this caused notorious failures (the PlayStation 3, early Bitcoin wallets).

ECDSA and its modern variant Ed25519 (using Edwards curves) are the dominant signature schemes for software signing, web certificates, SSH keys, and cryptocurrency wallets.

%------------------------------------------------------------------
\section{Curve choices and standards}
%------------------------------------------------------------------

Not all elliptic curves are equally secure or efficient. Standardised curves include:

\begin{itemize}
  \item \textbf{NIST P-256 (secp256r1).} A 256-bit prime-field curve published by NIST. Widely deployed in TLS, EMV, government systems. Suspect to some cryptographers for opaque parameter generation; the NSA's role in its design has been debated.
  \item \textbf{secp256k1.} The curve used by Bitcoin and many cryptocurrencies. Chosen for performance reasons (efficient endomorphism). Mathematically sound but not in original NIST suite.
  \item \textbf{Curve25519 / Ed25519.} A modern curve designed by Daniel Bernstein, providing high performance, ``rigid'' parameter generation, and resistance to side-channel attacks. Increasingly the default in new protocols.
  \item \textbf{Brainpool curves.} European-standardised curves with transparent parameter generation. Used in some European financial PKIs.
\end{itemize}

For finance applications, NIST P-256 dominates legacy systems (EMV cards, financial PKIs), with Ed25519 increasingly the choice for new protocols. The migration to post-quantum cryptography will affect both equally---no elliptic curve, however carefully chosen, resists Shor's algorithm.

%------------------------------------------------------------------
\section{Why ECC is the modern default}
%------------------------------------------------------------------

Compared to RSA, ECC offers concrete advantages at every security level:

\begin{itemize}
  \item \textbf{Smaller keys.} 256-bit ECC matches the security of 3072-bit RSA. For 192-bit security, ECC uses 384-bit keys vs RSA's 7680-bit keys.
  \item \textbf{Faster operations.} ECC point multiplication is much faster than RSA modular exponentiation at equivalent security.
  \item \textbf{Smaller signatures and ciphertexts.} An ECDSA signature is around 64 bytes; an RSA signature at equivalent security is 256+ bytes.
  \item \textbf{Lower power consumption.} Critical for smart cards, IoT, and mobile devices.
\end{itemize}

These advantages have driven ECC's dominance in new deployments since around 2010. TLS 1.3 essentially defaults to ECC for key exchange. Mobile banking apps use ECC for authentication. Bitcoin and most other cryptocurrencies use ECDSA on secp256k1. Smart cards in modern EMV chip and mobile-wallet deployments are ECC-based.

\begin{remark}
The story of ECC adoption in financial systems is a slow tide rather than a sudden flip. EMV chip cards transitioned from RSA-only to ECC-supporting over more than a decade. Many legacy systems still use RSA in parallel. The PQC migration will be an equally long process---perhaps another decade---and the financial industry will likely run all three (RSA, ECC, PQC) in parallel for years.
\end{remark}

%------------------------------------------------------------------
\section{Pairings and identity-based cryptography}
%------------------------------------------------------------------

A more exotic family of elliptic curves admit \keyterm{bilinear pairings}, mappings $e: E_1 \times E_2 \to \mathbb{G}_T$ satisfying $e([a]P, [b]Q) = e(P,Q)^{ab}$. Pairings enable cryptographic primitives unavailable with vanilla ECC:

\begin{itemize}
  \item \textbf{Identity-based encryption} (IBE), where a user's identity (e.g., email address) serves directly as their public key.
  \item \textbf{Short signatures} (BLS), used by Ethereum and other blockchain protocols.
  \item \textbf{Aggregate signatures}, where many signatures can be combined into one short signature.
\end{itemize}

Pairing-based schemes have been used in some financial blockchain protocols and in research-stage privacy systems. They too rely on the hardness of the ECDLP (or the related Bilinear Diffie--Hellman Problem) and are vulnerable to the same quantum attacks.

%------------------------------------------------------------------
\section{Shor's algorithm hits ECC at least as hard as RSA}
%------------------------------------------------------------------

Shor's algorithm (Chapter on Shor) is most famous for factoring, but it solves the discrete logarithm problem in any finite abelian group in polynomial time. The ECDLP is a discrete log in the group of points on an elliptic curve over $\mathbb{F}_p$, so Shor breaks it.

The resource estimates for breaking ECC with Shor are, in fact, smaller than for breaking RSA at equivalent classical security. A 256-bit ECC key (128-bit classical security) requires roughly 2500 logical qubits and $\sim 10^{10}$ Toffoli gates under recent optimised resource estimates---less than the $\sim 6000$ logical qubits and $\sim 10^{11}$ Toffoli gates for 2048-bit RSA. So ECC is, paradoxically, slightly more vulnerable to early fault-tolerant quantum computers than RSA at equivalent classical security.

This has practical consequences. An adversary with a moderately-sized fault-tolerant quantum computer would target ECC-secured systems first. Since ECC is now the dominant key-exchange algorithm in TLS, mobile banking, and cryptocurrencies, the financial sector's quantum exposure is concentrated in ECC just as much as in RSA.

\begin{example}
A Bitcoin transaction is signed with ECDSA on secp256k1. Once a transaction's public key is revealed on the blockchain (which happens at the moment of spending), a sufficiently powerful quantum computer could recover the private key and forge subsequent transactions. The community is actively researching post-quantum replacements (lattice-based and hash-based signature schemes) and migration paths.
\end{example}

%------------------------------------------------------------------
\section{The post-quantum landscape for ECC replacements}
%------------------------------------------------------------------

The same NIST post-quantum standards (ML-KEM, ML-DSA, SLH-DSA, FALCON) that replace RSA also replace ECC. For finance, the migration considerations are similar: inventory dependencies, prioritize long-confidentiality data, deploy hybrid protocols during transition, eventually retire ECC in production.

One difference matters. ECC has smaller key sizes and signatures than RSA. The post-quantum replacements typically have \emph{larger} keys and signatures than ECC---ML-KEM-768 has 1184-byte public keys, ML-DSA-65 has 1952-byte public keys, far above the 32--64 byte ECC analogues. For bandwidth-constrained protocols (mobile banking, smart cards), this presents engineering challenges: protocols must be updated to accommodate larger keys and signatures, and embedded systems may need hardware upgrades.

The migration from ECC to PQC will therefore be more disruptive in some sectors than the migration from RSA, despite ECC's modern reputation. Financial systems that already use ECC for performance reasons (smart cards, mobile devices) will feel this cost most.

%------------------------------------------------------------------
\section{The unifying lesson}
%------------------------------------------------------------------

RSA and ECC are mathematically different but share a common structural weakness: both rely on the difficulty of an abelian-group discrete logarithm or factoring problem, and both fall to Shor's algorithm. The next chapter explains how.

The broader cryptographic lesson is that \emph{algebraic structure is dangerous}. Any cryptosystem whose security can be reduced to a hidden-subgroup problem over an abelian group is vulnerable to Shor. Post-quantum cryptosystems must therefore base their hardness on different mathematical problems---lattice problems, code-based problems, hash-based constructions, multivariate polynomial systems---that resist the quantum techniques of Shor and his successors. The diversity of the NIST PQC finalists reflects this strategic shift: not putting all post-quantum eggs in one mathematical basket.

\textsep
